\subsection{Source-population quadrature and detector errors}
\label{sec:source-population-quadrature}

Cell masses connect a finite source population to an integral on its metric
carrier. A covering radius controls nearest-point geometry; an assignment
radius controls quadrature for a specified measure. These radii can differ.
For example, a golden-orbit point need not lie in its equal-volume permutation
cell. The estimates below use the actual cell assignment and its measure.

Let $\Omega$ be a metric window with a finite positive Borel measure $\mu$.
Choose a finite measurable partition $C_i$, representatives $s_i\in\Omega$,
and masses $v_i=\mu(C_i)$. A measurable convention assigns shared boundaries;
null modifications do not affect the statements. Representatives need not
belong to their cells. Let $f$ take values in a real Banach space and be
integrable on every cell. Set
\[
 Qf=\sum_i v_i f(s_i),\qquad
 D_i=\int_{C_i}d(s_i,x)\,d\mu(x).
\]

\begin{proposition}[Integrated cell-transport estimate]
\label{prop:source-population-cell-transport}
Suppose
$\|f(s_i)-f(x)\|\le K_i d(s_i,x)$ for almost every $x\in C_i$,
where $K_i\ge0$ and the displacement is integrable. Then
\begin{equation}
 \left\|Qf-\int_\Omega f\,d\mu\right\|
 \le\sum_i K_i D_i.
 \label{eq:source-population-first-moment}
\end{equation}
If $d(s_i,x)\le H_i$ almost everywhere on $C_i$, the right-hand side is
at most $\sum_i v_iK_iH_i$. In particular, for a globally $K$-Lipschitz
function and a uniform assignment bound $H$,
\begin{equation}
 \left\|Qf-\int_\Omega f\,d\mu\right\|
 \le \mu(\Omega)KH.
 \label{eq:source-population-uniform-quadrature}
\end{equation}
For a sequence of such partitions of the same measured window, $H_n\to0$
therefore gives $Q_nf\to\int_\Omega f\,d\mu$ for each integrable Lipschitz
test function. Neither nesting nor a fixed number of representatives is
required.
\end{proposition}
\begin{proof}
The difference equals
$\sum_i\int_{C_i}[f(s_i)-f(x)]\,d\mu(x)$.
The norm of a Bochner integral is bounded by the integral of the norm.
Apply the local Lipschitz bound and sum the resulting inequalities.
The uniform estimate uses $D_i\le v_iH_i$ and
$\sum_i v_i=\mu(\Omega)$. Its right-hand side tends to zero under the
stated refinement. The first-moment coefficient is sharp: on $[0,H]$,
with Lebesgue measure, one representative at zero and $f(x)=x$, the error
and displacement integral both equal $H^2/2$.
\end{proof}

The Lean module \texttt{SourcePopulationQuadrature} proves the actual
Bochner-integral estimates for finite cell measures, their specialization
to measurable disjoint partitions, and convergence under a vanishing
assignment bound. It also checks that a constant detector exposes any
error in total assigned mass. The source coordinates, partition geometry
and their displacement bounds must be established for each population.
These hypotheses are not replaced by an assumed quadrature limit.

For the golden permutation cells of side $L/q$, the established assignment
bound $H_q=2\sqrt3L/q$ and equal masses $L^3/q^3$ instantiate
\eqref{eq:source-population-uniform-quadrature}. Product time cells of
length $\Delta_q$ contribute the independent temporal error
$TL^3\Delta_q\operatorname{Lip}_t(f)$ on an aligned window of duration
$T$. The spatial term is $TL^3H_q\operatorname{Lip}_x(f)$.
An unmatched endpoint layer contributes at most
$L^3\Delta_q\|f\|_\infty$. This application retains the declared event
population and time-cell convention; extra routing or bookkeeping events
cannot be inserted into its density without another count estimate.
The module \texttt{GoldenSourceAssignment} proves the Fibonacci
approximation and floor agreement, constructs the residue-permuted
rectangular partition, and derives its equal Lebesgue masses and Euclidean
assignment bound. For each fixed globally Lipschitz Banach-valued detector,
it proves the actual golden-population quadrature estimate and its
refinement limit; the closed and half-open cube integrals agree almost
everywhere. The spacetime null-boundary interval and strict-pair count
specializations are analytic arguments.

\begin{proposition}[Localized radial-kernel quadrature]
\label{prop:source-population-local-kernel}
For $\epsilon>0$, write
\[
 w_x(y)=(1-d(x,y)/\epsilon)_+,
 \qquad g_x(y)=w_x(y)[f(x)-f(y)].
\]
If $f$ is a real $K$-Lipschitz function, then $g_x$ is $K$-Lipschitz
and vanishes outside the closed radius-$\epsilon$ ball about $x$.
With an assignment map $P(y)=s_i$ on $C_i$ satisfying
$d(P(y),y)\le H$, where $H\ge0$, its quadrature error obeys
\begin{equation}
 \left|\sum_i v_i g_x(s_i)-\int_\Omega g_x(y)\,d\mu(y)\right|
 \le K\!\int_{\Omega\cap B(x,\epsilon+H)}
                      d(P(y),y)\,d\mu(y)
 \le KH\,\mu(\Omega\cap B(x,\epsilon+H)).
 \label{eq:source-population-local-kernel}
\end{equation}
The same statement holds with closed balls; their boundaries are irrelevant
for Euclidean volume.
\end{proposition}
\begin{proof}
The tent weight lies in $[0,1]$ and satisfies
$|w_x(y)-w_x(z)|\le d(y,z)/\epsilon$. Order the two weights so that
$w_x(y)\le w_x(z)$. If both vanish the claim is immediate; otherwise
$d(x,z)<\epsilon$. The identity
\[
 g_x(y)-g_x(z)
 =w_x(y)[f(z)-f(y)]
  +[w_x(y)-w_x(z)][f(x)-f(z)]
\]
is bounded by
$Kd(y,z)[w_x(y)+d(x,z)/\epsilon]$. The bracket is at most one because
$w_x(y)\le w_x(z)=1-d(x,z)/\epsilon$. The same proof with $y,z$
exchanged covers the other order. If $d(x,y)\ge\epsilon+H$, both kernel terms
$g_x(y)$ and $g_x(P(y))$ vanish. Integrate the pointwise bound over the
remaining region. The tent Lipschitz estimates have Lean proofs; the
displayed support-restricted integral applies the cell-transport argument.
\end{proof}

For three-dimensional volume, multiplying
\eqref{eq:source-population-local-kernel} by $45/(\pi\epsilon^5)$ gives
the quadrature contribution
$60KH(\epsilon+H)^3/\epsilon^5$ in the scalar operator estimate,
and hence $480KH/\epsilon^2$ when $H\le\epsilon$.
Actual local cell masses or integrated displacements can give smaller
bounds than the enclosing-ball volume. They must be computed or bounded
for the same assignment. A small fill radius does not justify replacing
the equal-volume assignment radius. This estimate alone yields no vanishing
operator error when $H=O(q^{-1})$ and $\epsilon=O(q^{-1/2})$.

\paragraph{Dual cells for a local tensor action.}
A different supplied action can use sorted source coordinates
$0=x_0<\cdots<x_N=L$, with interval lengths $\ell_j=x_{j+1}-x_j$.
Half-interval endpoint weights and interior weights
$(\ell_{j-1}+\ell_j)/2$ are the masses of the one-dimensional dual cells.
Their tensor products satisfy the preceding transport theorem. For
$f\in C^2([0,L])$, linear interpolation on each interval also gives
the sharper analytic trapezoidal bound
\[
 \left|Qf-\int_0^L f(x)\,dx\right|
 \le\frac{\|f''\|_\infty}{12}\sum_j\ell_j^3
 \le\frac{L\ell_{\max}^2}{12}\|f''\|_\infty.
\]
Indeed the pointwise interpolation error is at most
$\|f''\|_\infty(x-x_j)(x_{j+1}-x)/2$; integrate it on each interval.
Successive application in three coordinates gives the tensor bound
$L^3\ell_{\max}^2\sum_{a=1}^3\|\partial_a^2f\|_\infty/12$.
This applies to detector products and sampled mode Gram entries with their
stated derivative bounds. Dual-cell weights and equal-volume source-count
weights are different finite measures, each with its own comparison to
volume. The trapezoidal argument does not improve the radial action by
changing its weights implicitly, and it supplies no source-selection law,
primitive timing rule or physical calibration.
